\BeleskeID{09_1-s053}
% element: 09_1-s053:n01
Porast ulaznog napona snižava napon na drejnu pogonskog tranzistora. Posle prikazane granice on radi u omskoj oblasti, a opterećenje ostaje u zasićenju. Ravnoteža struja pre zamene napona glasi
\[\frac{k_{n,L}}2(V_{GS,L}-V_{Tn,L})^2=\frac{k_{n,D}}2\bigl(2V_{DS,D}(V_{GS,D}-V_{Tn,D})-V_{DS,D}^2\bigr)\]
Na levoj strani je kvadrat napona prezasićenja opterećenja; na desnoj je omski model pogonskog tranzistora. Pri diferenciranju se na desnoj strani koristi pravilo za izvod proizvoda. Druga relacija služi za određivanje radne tačke u kojoj je nagib $-1$.

\[2k_{n,L}(V_{DD}-V_O-V_{Tn,L})=-k_{n,D}(2V_I-2V_{Tn,L}-V_O)+k_{n,D}V_O(2+1)\]

